\documentclass[a4paper,11pt]{article}
\usepackage{amsmath,amsthm,amsfonts,amssymb,amscd,amstext,vmargin,graphics,graphicx,tabularx,multicol} 
\usepackage[francais]{babel}
\usepackage[utf8]{inputenc}  
\usepackage[T1]{fontenc} 
\usepackage{pstricks-add,tikz,tkz-tab,variations}
\usepackage[autolanguage,np]{numprint} 
\usepackage{calc}

\setmarginsrb{1.5cm}{0.5cm}{1cm}{0.5cm}{0cm}{0cm}{0cm}{0cm} %Gauche, haut, droite, haut
\newcounter{numexo}
\newcommand{\exo}[1]{\stepcounter{numexo}\noindent{\bf Exercice~\thenumexo} : }
\reversemarginpar

\newcommand{\bmul}[1]{\begin{multicols}{#1}}
\newcommand{\emul}{\end{multicols}}

\newcounter{enumtabi}
\newcounter{enumtaba}
\newcommand{\q}{\stepcounter{enumtabi} \theenumtabi.  }
\newcommand{\qa}{\stepcounter{enumtaba} (\alph{enumtaba}) }
\newcommand{\initq}{\setcounter{enumtabi}{0}}
\newcommand{\initqa}{\setcounter{enumtaba}{0}}

\newcommand{\be}{\begin{enumerate}}
\newcommand{\ee}{\end{enumerate}}
\newcommand{\bi}{\begin{itemize}}
\newcommand{\ei}{\end{itemize}}
\newcommand{\bp}{\begin{pspicture*}}
\newcommand{\ep}{\end{pspicture*}}
\newcommand{\bt}{\begin{tabular}}
\newcommand{\et}{\end{tabular}}
\renewcommand{\tabularxcolumn}[1]{>{\centering}m{#1}} %(colonne m{} centrée, au lieu de p par défault) 
\newcommand{\tnl}{\tabularnewline}

\newcommand{\trait}{\noindent \rule{\linewidth}{0.2mm}}
\newcommand{\hs}[1]{\hspace{#1}}
\newcommand{\vs}[1]{\vspace{#1}}

\newcommand{\N}{\mathbb{N}}
\newcommand{\Z}{\mathbb{Z}}
\newcommand{\R}{\mathbb{R}}
\newcommand{\C}{\mathbb{C}}
\newcommand{\Dcal}{\mathcal{D}}
\newcommand{\Ccal}{\mathcal{C}}
\newcommand{\mc}{\mathcal}

\newcommand{\vect}[1]{\overrightarrow{#1}}
\newcommand{\ds}{\displaystyle}
\newcommand{\eq}{\quad \Leftrightarrow \quad}
\newcommand{\vecti}{\vec{\imath}}
\newcommand{\vectj}{\vec{\jmath}}
\newcommand{\Oij}{(O;\vec{\imath}, \vec{\jmath})}
\newcommand{\OIJ}{(O;I,J)}


\newcommand{\reponse}[1][1]{%
\multido{}{#1}{\makebox[\linewidth]{\rule[0pt]{0pt}{20pt}\dotfill}
}}

\newcommand{\titre}[5] 
% #1: titre #2: haut gauche #3: bas gauche #4: haut droite #5: bas droite
{
\noindent #2 \hfill #4 \\
#3 \hfill #5

\vspace{-1.6cm}

\begin{center}\rule{6cm}{0.5mm}\end{center}
\vspace{0.2cm}
\begin{center}{\large{\textbf{#1}}}\end{center}
\begin{center}\rule{6cm}{0.5mm}\end{center}
}



\begin{document}
\pagestyle{empty}
\titre{Séance d'AP 4 : Transformations du plan}{}{}{4ème}{}

\vspace*{0.2cm}

\setlength{\fboxrule}{2pt}
\begin{flushleft}
\framebox{\begin{minipage}{\linewidth}

\vspace*{0.2cm}

\underline{\textbf{Rappels de cours}}\\

\underline{Définition} : On appelle polygone régulier, un polygone dont :
\bi \item tous les côtés ont la même longueur,
\item tous les angles ont la même mesure.\\
\ei

\q  Quels sont les polygones réguliers à 3 côtés ? . . . . . . . . . . . . . . . . . . . . . . .\\
\q  Quels sont les polygones réguliers à 4 côtés ? . . . . . . . . . . . . . . . . . . . . . . .\\
\q  Quels sont les polygones réguliers à 5 côtés ? . . . . . . . . . . . . . . . . . . . . . . .\\
\vspace*{0.3cm}

\underline{Propriété :} Un polygone est régulier si, et seulement s'il existe . . . . . . . . . . . . . . . . . . . . . . . . . . . . qui envoie chaque sommet sur le suivant.\\



\end{minipage}}
\end{flushleft}

\vspace*{0.4cm}

OUVRIR LE LOGICIEL GEOGEBRA.\includegraphics[scale=0.5]{geogebra.eps} \\
Ce logiciel va nous permettre de tracer des figures géométriques et de les « faire bouger ».\\


\exo \\

\initq 
\noindent \q \qa  Tracer ABCD un carré  de 4 cm de côté.\\
\qa Tracer  4 triangles équilatéraux DCE, BCF, ABG, ADH.

\hspace*{5cm} \includegraphics[scale=1]{notetp.eps} 

\noindent \qa Tracer la parallèle à (BC) passant par F. Puis la parallèle (AB) passant par G. De la même manière, la parallèle à (AD) passant par H et enfin la parallèle à (DC) passant par E.\\
\qa Pour la suite de l'exercice, noter les points d'intersection des parallèles I, J, K et L.\\

\hspace*{5cm} \includegraphics[scale=1]{notetp.eps} 

\noindent \q \initqa \qa Construire les images de la figure précédemment tracée par la symétrie d'axe (JK), (LK), (IL), (IJ).\\

\hspace*{5cm} \includegraphics[scale=1]{notetp.eps} 

\noindent \qa Puis construire les images de la figure par la symétrie de centres I, J, K et L.\\
\qa Vous venez de réaliser \textbf{un pavage}, c'est-à-dire,  le remplissage d'un plan avec un motif qui se répète. \\
Vous pouvez agrandir ce pavage en répétant indéfiniment les étapes que vous venez de faire et le coloriez suivant vos goûts.\\

\hspace*{5cm} \includegraphics[scale=1]{notetp.eps} \\

\exo\\

\initq 
\noindent \q \initqa \qa Construire un triangle ABC équilatéral de 3 cm de côté.\\
\qa Construire l'image du triangle ABC par la translation qui envoie le point A sur le point B .\\

\hspace*{5cm} \includegraphics[scale=1]{notetp.eps} 
\newpage

\vspace*{0.5cm}

\noindent \qa Construire l'image du triangle ABC par la rotation de centre B et d'angle 120\degre dans le sens anti-horaire.\\

\hspace*{5cm} \includegraphics[scale=1]{notetp.eps} 

\noindent \q \initqa \qa A l'aide des points déjà placés construire un hexagone.\\
 \qa Prouver que les longueurs des côtés de l'hexagone sont toutes égales.\\
 \reponse[5]\\
 
\noindent \qa A l'aide de transformations bien choisies, construire un pavage du plan à l'aide d'hexagones (colorer à votre convenance).\\
 
 \hspace*{5cm} \includegraphics[scale=1]{notetp.eps} 



\end{document}
